\subsection{RLC Band Pass Filter}
\label{sec:rlc_band_pass}

\paragraph{Overview:}
Constructed from a resistor, inductor, and capacitor in series
or parallel, a RLC Band Pass Filter allows a specific band of
frequencies to pass while attenuating frequencies outside this
band. It is a second order filter.

\paragraph{Schematic:}
\begin{center}
  \begin{circuitikz}
    % Paths, nodes and wires:
	  \draw (7.5, 5.75) to[american voltage source, l_={$V_{ee}$}] (6.25, 5.75);
	  \node[ground] at (4.75, 6.75){};
	  \node[ground, rotate=-90] at (6.25, 5.75){};
	  \draw (4.75, 8) -| (4.75, 8.25);
	  \draw (8.25, 8.25) -- (9, 8.25);
	  \draw (4.75, 6.75) to[sinusoidal voltage source, l_={$V_{in}$}] (4.75, 8);
	  \node[shape=rectangle, minimum width=0.965cm, minimum height=0.465cm](N1) at (9.5, 8.25){} node[anchor=center] at (N1.text){$V_{out}$};
	  \node[shape=rectangle, minimum width=3.215cm, minimum height=0.465cm] at (6.75, 10.75){} node[anchor=north, align=center, text width=2.827cm, inner sep=6pt] at (6.75, 11){Full Schematic};
	  \node[shape=rectangle, minimum width=5.215cm, minimum height=0.465cm] at (13.375, 10.75){} node[anchor=north, align=center, text width=4.827cm, inner sep=6pt] at (13.375, 11){Small-Signal Schematic};
	  \draw (7.75, 6.25) to[american inductor, l={$L_1$}] (7.75, 7.5);
	  \draw (8.75, 6.25) to[capacitor, l_={$C_1$}] (8.75, 7.5);
	  \draw (8.75, 6.25) -- (8.75, 6) |- (8.25, 6);
	  \draw (8.25, 6) |- (7.75, 6) -- (7.75, 6.25);
	  \draw (7.75, 7.5) -- (7.75, 7.75) |- (8.25, 7.75);
	  \draw (8.75, 7.5) |- (8.25, 7.75);
	  \draw (4.75, 8.25) -- (5.75, 8.25);
	  \draw (7.5, 8.25) to[european resistor, l={$G_1$}] (5.75, 8.25);
	  \draw (7.5, 5.75) -| (8.25, 6);
	  \draw (8.25, 7.75) -| (8.25, 8.25) -- (7.5, 8.25);
	  \draw (7.5, 1.25) to[american voltage source, l_={$V_{ee}$}] (6.25, 1.25);
	  \node[ground] at (4.75, 2.25){};
	  \node[ground, rotate=-90] at (6.25, 1.25){};
	  \draw (4.75, 3.5) -| (4.75, 3.75);
	  \draw (8.25, 3.75) -- (9, 3.75);
	  \draw (4.75, 2.25) to[sinusoidal voltage source, l_={$V_{in}$}] (4.75, 3.5);
	  \node[shape=rectangle, minimum width=0.965cm, minimum height=0.465cm](N2) at (9.5, 3.75){} node[anchor=center] at (N2.text){$V_{out}$};
	  \draw (5.25, 3.75) to[american inductor, l={$L_1$}] (6.5, 3.75);
	  \draw (6.5, 3.75) to[capacitor, l={$C_1$}] (7.75, 3.75);
	  \draw (4.75, 3.75) -- (5.25, 3.75);
	  \draw (8.25, 3.25) to[european resistor, l={$G_1$}] (8.25, 1.5);
	  \draw (7.5, 1.25) -| (8.25, 1.5);
	  \draw (8.25, 3.25) -| (8.25, 3.75) -- (7.75, 3.75);
	  \draw (6.5, 3.75) -| (6.5, 3);
	  \node[shape=rectangle, minimum width=0.965cm, minimum height=0.465cm](N3) at (6.5, 2.75){} node[anchor=center] at (N3.text){$V_{x}$};
	  \node[shape=rectangle, minimum width=3.215cm, minimum height=0.465cm] at (2.125, 7.5){} node[anchor=north, align=center, text width=2.827cm, inner sep=6pt] at (2.125, 7.75){Parallel};
	  \node[shape=rectangle, minimum width=3.215cm, minimum height=0.465cm] at (2.25, 3.25){} node[anchor=north, align=center, text width=2.827cm, inner sep=6pt] at (2.25, 3.5){Series};
	  \draw (10, 11) -- (10, 0.5);
	  \draw (1, 5) -- (17, 5);
	  \draw (4, 0.5) -- (4, 11);
	  \draw (1, 10) -- (17, 10);
	  \node[ground] at (10.75, 6.875){};
	  \node[ground, rotate=-90] at (13.5, 5.875){};
	  \draw (10.75, 8.125) -| (10.75, 8.375);
	  \draw (14.25, 8.375) -- (15, 8.375);
	  \draw (10.75, 6.875) to[sinusoidal voltage source, l_={$V_{in}$}] (10.75, 8.125);
	  \node[shape=rectangle, minimum width=0.965cm, minimum height=0.465cm](N4) at (15.5, 8.375){} node[anchor=center] at (N4.text){$V_{out}$};
	  \draw (13.75, 6.375) to[american inductor, l={$L_1$}] (13.75, 7.625);
	  \draw (14.75, 6.375) to[capacitor, l_={$C_1$}] (14.75, 7.625);
	  \draw (14.75, 6.375) -- (14.75, 6.125) |- (14.25, 6.125);
	  \draw (14.25, 6.125) |- (13.75, 6.125) -- (13.75, 6.375);
	  \draw (13.75, 7.625) -- (13.75, 7.875) |- (14.25, 7.875);
	  \draw (14.75, 7.625) |- (14.25, 7.875);
	  \draw (10.75, 8.375) -- (11.75, 8.375);
	  \draw (13.5, 8.375) to[european resistor, l={$G_1$}] (11.75, 8.375);
	  \draw (13.5, 5.875) -| (14.25, 6.125);
	  \draw (14.25, 7.875) -| (14.25, 8.375) -- (13.5, 8.375);
	  \node[ground] at (10.75, 2.25){};
	  \node[ground, rotate=-90] at (13.5, 1.25){};
	  \draw (10.75, 3.5) -| (10.75, 3.75);
	  \draw (14.25, 3.75) -- (15, 3.75);
	  \draw (10.75, 2.25) to[sinusoidal voltage source, l_={$V_{in}$}] (10.75, 3.5);
	  \node[shape=rectangle, minimum width=0.965cm, minimum height=0.465cm](N5) at (15.5, 3.75){} node[anchor=center] at (N5.text){$V_{out}$};
	  \draw (11.25, 3.75) to[american inductor, l={$L_1$}] (12.5, 3.75);
	  \draw (12.5, 3.75) to[capacitor, l={$C_1$}] (13.75, 3.75);
	  \draw (10.75, 3.75) -- (11.25, 3.75);
	  \draw (14.25, 3.25) to[european resistor, l={$G_1$}] (14.25, 1.5);
	  \draw (13.5, 1.25) -| (14.25, 1.5);
	  \draw (14.25, 3.25) -| (14.25, 3.75) -- (13.75, 3.75);
	  \draw (12.5, 3.75) -| (12.5, 3);
	  \node[shape=rectangle, minimum width=0.965cm, minimum height=0.465cm](N6) at (12.5, 2.75){} node[anchor=center] at (N6.text){$V_{x}$};
  \end{circuitikz}
\end{center}

\subsubsection{Transfer Function}
\begin{enumerate}
  \item \textbf{Parallel Configuration:}
    Applying KCL at $V_{out}$ on Full Schematic of Parallel
    Configuration, we get:
    \begin{align*}
      (V_{out} - V_{in}){G_1} + \frac{(V_{out} - V_{ee})}{sL_1} + (V_{out} - V_{ee}){sC_1} &= 0 \\
      V_{out}G_1 - V_{in}G_1 + \frac{V_{out}}{sL_1} - \frac{V_{ee}}{sL_1} + V_{out}sC_1 - V_{ee}sC_1 &= 0 \\
    \end{align*}
    \paragraph{Superposition:}
      \begin{enumerate}
        \item Due to \(V_{in}\) alone (\(V_{ee} = 0V\)):
          \begin{align*}
            V_{out}G_1 - V_{in}G_1 + \frac{V_{out}}{sL_1} + V_{out}sC_1 &= 0 \\
            V_{out}(G_1 + \frac{1}{sL_1} + sC_1) - V_{in}G_1 &= 0 \\
            V_{out}(G_1 + \frac{1}{sL_1} + sC_1) &= V_{in}G_1 \\
            \boxed{H_1(s) = \frac{N_1(s)}{D_1(s)} = \frac{V_{out_1}}{V_{in}} = \frac{G_1}{G_1 + \frac{1}{sL_1} + sC_1}}
          \end{align*}
        \item Due to \(V_{ee}\) alone (\(V_{in} = 0V\)):
          \begin{align*}
            V_{out}G_1 + \frac{V_{out}}{sL_1} - \frac{V_{ee}}{sL_1} + V_{out}sC_1 - V_{ee}sC_1 &= 0 \\
            V_{out}(G_1 + \frac{1}{sL_1} + sC_1) - V_{ee}(\frac{1}{sL_1} + sC_1) &= 0 \\
            V_{out}(G_1 + \frac{1}{sL_1} + sC_1) &= V_{ee}(\frac{1}{sL_1} + sC_1) \\
            \boxed{H_2(s) = \frac{N_2(s)}{D_2(s)} = \frac{V_{out_2}}{V_{ee}} = \frac{\frac{1}{sL_1} + sC_1}{G_1 + \frac{1}{sL_1} + sC_1}}
          \end{align*}
      \end{enumerate}
    \paragraph{Total Parallel Response:}
      \begin{align*}
        V_{out} &= V_{out_1} + V_{out_2} \\
                &= H_1(s)V_{in} + H_2(s)V_{ee} \\
                &= \left(\frac{G_1}{G_1 + \frac{1}{sL_1} + sC_1}\right)V_{in} + \left(\frac{\frac{1}{sL_1} + sC_1}{G_1 + \frac{1}{sL_1} + sC_1}\right)V_{ee} \\
        \boxed{V_{out} = \frac{G_1}{G_1 + \frac{1}{sL_1} + sC_1}V_{in} + \frac{\frac{1}{sL_1} + sC_1}{G_1 + \frac{1}{sL_1} + sC_1}V_{ee}}
      \end{align*}
  \item \textbf{Series Configuration:}
    Applying KCL at $V_{out}$ on Full Schematic of Series
    Configuration, we get:
    \begin{align*}
      (V_{out} - V_{x})sC_1 + (V_{out} - V_{ee}){G_1} &= 0 \tag{i}
    \end{align*}
    Applying KCL at \(V_{x}\):
    \begin{align*}
      \frac{(V_{x} - V_{in})}{sL_1} + (V_{x} - V_{out})sC_1 &= 0 \tag{ii}
    \end{align*}
    From equation (ii):
    \begin{align*}
      \frac{V_{x}}{sL_1} - \frac{V_{in}}{sL_1} + V_{x}sC_1 - V_{out}sC_1 &= 0 \\
      V_{x}(\frac{1}{sL_1} + sC_1) - \frac{V_{in}}{sL_1} - V_{out}sC_1 &= 0 \\
      V_{x}(\frac{1}{sL_1} + sC_1) &= \frac{V_{in}}{sL_1} + V_{out}sC_1 \\
      V_{x} &= \frac{\frac{V_{in}}{sL_1} + V_{out}sC_1}{\frac{1}{sL_1} + sC_1} \\
      V_{x} &= \frac{\left(\frac{V_{in} + V_{out}s^2 L_1 C_1}{s L_1}\right)}{\left(\frac{1 + s^2 L_1 C_1}{s L_1}\right)} \\
      V_{x} &= \frac{V_{in} + V_{out}s^2 L_1 C_1}{1 + s^2 L_1 C_1} \tag{iii}
    \end{align*}
    \paragraph{Superposition:}
      \begin{enumerate}
        \item Due to \(V_{in}\) alone (\(V_{ee} = 0V\)):
          Substituting equation (iii) into equation (i):
          \begin{align*}
            \left(V_{out} - \left(\frac{V_{in} + V_{out}s^2 L_1 C_1}{1 + s^2 L_1 C_1}\right)\right) \cdot sC_1 + (V_{out} - 0){G_1} &= 0 \\
            \left(V_{out}(1 + s^2 L_1 C_1) - (V_{in} + V_{out}s^2 L_1 C_1)\right) \cdot \frac{sC_1}{1 + s^2 L_1 C_1} + V_{out}G_1 &= 0 \\
            \left(V_{out} + V_{out}s^2 L_1 C_1 - V_{in} - V_{out}s^2 L_1 C_1\right) \cdot \frac{sC_1}{1 + s^2 L_1 C_1} + V_{out}G_1 &= 0 \\
            \left(V_{out} - V_{in}\right) \cdot \frac{sC_1}{1 + s^2 L_1 C_1} + V_{out}G_1 &= 0 \\
            \frac{V_{out} sC_1}{1 + s^2 L_1 C_1} - \frac{V_{in} sC_1}{1 + s^2 L_1 C_1} + V_{out}G_1 &= 0 \\
            V_{out}\left(\frac{sC_1}{1 + s^2 L_1 C_1} + G_1\right) &= \frac{V_{in} sC_1}{1 + s^2 L_1 C_1} \\
            V_{out}\left(\frac{sC_1 + G_1(1 + s^2 L_1 C_1)}{1 + s^2 L_1 C_1}\right) &= \frac{V_{in} sC_1}{1 + s^2 L_1 C_1} \\
            V_{out} &= \frac{\frac{V_{in} sC_1}{1 + s^2 L_1 C_1}}{\frac{sC_1 + G_1(1 + s^2 L_1 C_1)}{1 + s^2 L_1 C_1}} \\
            V_{out} = \frac{V_{in} sC_1}{sC_1 + G_1(1 + s^2 L_1 C_1)}
          \end{align*}
          This results in:
          \[
            \boxed{H_3(s) = \frac{N_3(s)}{D_3(s)} = \frac{V_{out_3}}{V_{in}} = \frac{sC_1}{sC_1 + G_1(1 + s^2 L_1 C_1)}}
          \]
        \item Due to \(V_{ee}\) alone (\(V_{in} = 0V\)):
          Substituting equation (iii) into equation (i):
          \begin{align*}
            \left(V_{out} - \left(\frac{0 + V_{out}s^2 L_1 C_1}{1 + s^2 L_1 C_1}\right)\right) \cdot sC_1 + (V_{out} - V_{ee}){G_1} &= 0 \\
            \left(V_{out}(1 + s^2 L_1 C_1) - V_{out}s^2 L_1 C_1\right) \cdot \frac{sC_1}{1 + s^2 L_1 C_1} + (V_{out} - V_{ee}){G_1} &= 0 \\
            \left(V_{out} + V_{out}s^2 L_1 C_1 - V_{out}s^2 L_1 C_1\right) \cdot \frac{sC_1}{1 + s^2 L_1 C_1} + (V_{out} - V_{ee}){G_1} &= 0 \\
            V_{out} \cdot \frac{sC_1}{1 + s^2 L_1 C_1} + (V_{out} - V_{ee}){G_1} &= 0 \\
            \frac{V_{out} sC_1}{1 + s^2 L_1 C_1} + V_{out}G_1 - V_{ee}G_1 &= 0 \\
            V_{out}\left(\frac{sC_1}{1 + s^2 L_1 C_1} + G_1\right) &= V_{ee}G_1 \\
            V_{out}\left(\frac{sC_1 + G_1(1 + s^2 L_1 C_1)}{1 + s^2 L_1 C_1}\right) &= V_{ee}G_1 \\
            V_{out} &= \frac{V_{ee}G_1}{\frac{sC_1 + G_1(1 + s^2 L_1 C_1)}{1 + s^2 L_1 C_1}} \\
            V_{out} = \frac{V_{ee}G_1(1 + s^2 L_1 C_1)}{sC_1 + G_1(1 + s^2 L_1 C_1)}
          \end{align*}
          This results in:
          \[
            \boxed{H_4(s) = \frac{N_4(s)}{D_4(s)} = \frac{V_{out_4}}{V_{ee}} = \frac{G_1(1 + s^2 L_1 C_1)}{sC_1 + G_1(1 + s^2 L_1 C_1)}}
          \]
      \end{enumerate}
    \paragraph{Total Series Response:}
      \begin{align*}
        V_{out} &= V_{out_3} + V_{out_4} \\
                &= H_3(s)V_{in} + H_4(s)V_{ee} \\
                = \left(\frac{sC_1}{sC_1 + G_1(1 + s^2 L_1 C_1)}\right)V_{in} + \left(\frac{G_1(1 + s^2 L_1 C_1)}{sC_1 + G_1(1 + s^2 L_1 C_1)}\right)V_{ee} \\
                \boxed{V_{out} = \frac{sC_1}{sC_1 + G_1(1 + s^2 L_1 C_1)}V_{in} + \frac{G_1(1 + s^2 L_1 C_1)}{sC_1 + G_1(1 + s^2 L_1 C_1)}V_{ee}}
      \end{align*}
\end{enumerate}

\subsubsection{Analysis}
\begin{enumerate}
  \item \textbf{Poles (\(D(s) = 0\)):} \\
      \begin{enumerate}
        \item \textbf{Parallel Configuration:}
          For both \(H_1(s)\) and \(H_2(s)\):
          \begin{align*}
            G_1 + \frac{1}{sL_1} + sC_1 &= 0 \\
            \frac{s L_1 G_1 + 1 + s^2 L_1 C_1}{s L_1} &= 0 \\
            s^2 L_1 C_1 + s L_1 G_1 + 1 &= 0 \\
            s^2 + s \frac{G_1}{C_1} + \frac{1}{L_1 C_1} &= 0 \tag{iv}
          \end{align*}
          Applying Shree Dharacharya's Formula:
          \begin{align*}
            s &= \frac{-\frac{G_1}{C_1} \pm \sqrt{\left(\frac{G_1}{C_1}\right)^2 - 4 \cdot 1 \cdot \frac{1}{L_1 C_1}}}{2 \cdot 1} \\
            s &= \frac{-\frac{G_1}{C_1} \pm \sqrt{\frac{G_1^2}{C_1^2} - \frac{4}{L_1 C_1}}}{2} \\
            \boxed{s = \frac{-G_1 \pm \sqrt{G_1^2 - \frac{4 C_1}{L_1}}}{2 C_1}}
          \end{align*}
          While the nature of the poles depends on the value of the
          discriminant \(\left(G_1^2 - \frac{4 C_1}{L_1}\right)\).
          We can still make a comparison of equation (iv) with natural
          second order system equation.\\ \\
          \begin{tabularx}{\linewidth}{|X|X|X|X|} \hline
            Type & Term A & Term B & Term C \\ \hline
            Equation (iv) & \(s^2\) & \(\frac{G_1}{C_1} s\) & \(\frac{1}{L_1 C_1}\) \\ \hline
            Natural & \(s^2\) & \(2 \zeta \omega_n s\) & \(\omega_n^2\) \\ \hline
          \end{tabularx}\\ \\
          We get:
          \begin{enumerate}
            \item Natural Frequency:
              \begin{align*}
                \omega_n = \sqrt{\frac{1}{L_1 C_1}} \\
                f_p = \frac{\omega_n}{2\pi} = \frac{1}{2\pi \sqrt{L_1 C_1}}
              \end{align*}
            \item Damping Ratio:
              \[
                2 \zeta \omega_n = \frac{G_1}{C_1} \implies \zeta = \frac{G_1}{2C_1 \omega_n} = \frac{G_1}{2 C_1} \sqrt{L_1 C_1}
              \]
            \item Q Factor:
              \[
                Q = \frac{1}{2\zeta} = \frac{C_1}{G_1} \sqrt{\frac{1}{L_1 C_1}}
              \]
          \end{enumerate}
        \item \textbf{Series Configuration:}
          For both \(H_3(s)\) and \(H_4(s)\):
          \begin{align*}
            sC_1 + G_1(1 + s^2 L_1 C_1) &= 0 \\
            G_1 s^2 L_1 C_1 + sC_1 + G_1 &= 0 \\
            s^2 L_1 C_1 G_1 + s C_1 + G_1 &= 0 \\
            s^2 + s \frac{1}{L_1 G_1} + \frac{1}{L_1 C_1} &= 0 \tag{v}
          \end{align*}
          Applying Shree Dharacharya's Formula:
          \begin{align*}
            s &= \frac{-\frac{1}{L_1 G_1} \pm \sqrt{\left(\frac{1}{L_1 G_1}\right)^2 - 4 \cdot 1 \cdot \frac{1}{L_1 C_1}}}{2 \cdot 1} \\
            s &= \frac{-\frac{1}{L_1 G_1} \pm \sqrt{\frac{1}{L_1^2 G_1^2} - \frac{4}{L_1 C_1}}}{2} \\
            s &= \frac{-\frac{1}{L_1 G_1} \pm \sqrt{\frac{1}{L_1^2 G_1^2} \left(1 - \frac{4 G_1^2 L_1}{C_1}\right)}}{2} \\
            \boxed{s = \frac{-1 \pm \sqrt{1 - \frac{4 G_1^2 L_1}{C_1}}}{2 L_1 G_1}}
          \end{align*}
          While the nature of the poles depends on the value of the
          discriminant \(\left(1 - \frac{4 G_1^2 L_1}{C_1}\right)\).
          We can still make a comparison of equation (v) with natural
          second order system equation.\\ \\
          \begin{tabularx}{\linewidth}{|X|X|X|X|} \hline
            Type & Term A & Term B & Term C \\ \hline
            Equation (v) & \(s^2\) & \(\frac{1}{L_1 G_1} s\) & \(\frac{1}{L_1 C_1}\) \\ \hline
            Natural & \(s^2\) & \(2 \zeta \omega_n s\) & \(\omega_n^2\) \\ \hline
          \end{tabularx}\\ \\
          We get,
          \begin{enumerate}
            \item Natural Frequency:
              \begin{align*}
                \omega_n = \sqrt{\frac{1}{L_1 C_1}} \\
                f_p = \frac{\omega_n}{2\pi} = \frac{1}{2\pi \sqrt{L_1 C_1}}
              \end{align*}
            \item Damping Ratio:
              \[
                2 \zeta \omega_n = \frac{1}{L_1 G_1} \implies \zeta = \frac{1}{2 L_1 G_1 \omega_n} = \frac{1}{2 G_1} \sqrt{\frac{C_1}{L_1}}
              \]
            \item Q Factor:
              \[
                Q = \frac{1}{2\zeta} = G_1 \sqrt{\frac{L_1}{C_1}}
              \]
          \end{enumerate}
      \end{enumerate}
    \item \textbf{Zeros (\(N(s) = 0\)):} \\
      \begin{enumerate}
        \item \textbf{Parallel Configuration:}\\
          For \(H_1(s)\):
            \begin{align*}
              G_1 &= 0 \\
            \end{align*}
            No zeros for \(H_1(s)\). \\
          For \(H_2(s)\):
            \begin{align*}
              \frac{1}{sL_1} + sC_1 &= 0 \\
              s^2 C_1 L_1 + 1 &= 0 \\
              s^2 + \frac{1}{L_1 C_1} &= 0 \\
              s &= \pm j\sqrt{\frac{1}{L_1 C_1}} \quad \text{[let \(s = s_z\)]} \\
              \boxed{f_z = \frac{|s_z|}{2\pi} = \frac{\sqrt{\frac{1}{L_1 C_1}}}{2\pi} = \frac{1}{2\pi \sqrt{L_1 C_1}}}
            \end{align*}
            Zero at natural frequency for \(H_2(s)\).
          \item \textbf{Series Configuration:}\\
          For \(H_3(s)\):
            \begin{align*}
              sC_1 &= 0 \\
              s &= 0
            \end{align*}
            Zero at origin for \(H_3(s)\). \\
          For \(H_4(s)\):
            \begin{align*}
              G_1(1 + s^2 L_1 C_1) &= 0 \\
              1 + s^2 L_1 C_1 &= 0 \\
              s^2 + \frac{1}{L_1 C_1} &= 0 \\
              s &= \pm j\sqrt{\frac{1}{L_1 C_1}} \quad \text{[let \(s = s_z\)]} \\
              \boxed{f_z = \frac{|s_z|}{2\pi} = \frac{\sqrt{\frac{1}{L_1 C_1}}}{2\pi} = \frac{1}{2\pi \sqrt{L_1 C_1}}}
            \end{align*}
           Zero at natural frequency for \(H_4(s)\).
      \end{enumerate}
    \item \textbf{DC Gain (\(H(0)\)):} \\
      \begin{enumerate}
      \item \textbf{Parallel Configuration:}\\
          For \(H_1(s)\):
          \begin{align*}
            H_1(0) &= \frac{G_1}{G_1 + \frac{1}{0 \cdot L_1} + 0 \cdot C_1} = \frac{G_1}{G_1 + \infty} = 0
          \end{align*}
          Zero DC gain for \(H_1(s)\). \\
          For \(H_2(s)\):
          \begin{align*}
            H_2(0) &= \frac{\frac{1}{0 \cdot L_1} + 0 \cdot C_1}{G_1 + \frac{1}{0 \cdot L_1} + 0 \cdot C_1} = \frac{\infty + 0}{G_1 + \infty} = 1
          \end{align*}
          Unity DC gain for \(H_2(s)\).
        \item \textbf{Series Configuration:}\\
          For \(H_3(s)\):
          \begin{align*}
            H_3(0) &= \frac{0 \cdot C_1}{0 \cdot C_1 + G_1(1 + 0^2 \cdot L_1 C_1)} = 0
          \end{align*}
          Zero DC gain for \(H_3(s)\). \\
          For \(H_4(s)\):
          \begin{align*}
            H_4(0) &= \frac{G_1(1 + 0^2 \cdot L_1 C_1)}{0 \cdot C_1 + G_1(1 + 0^2 \cdot L_1 C_1)} = \frac{G_1(1 + 0)}{0 + G_1(1 + 0)} = 1 \\
          \end{align*}
          Unity DC gain for \(H_4(s)\).
      \end{enumerate}
    \item \textbf{Gain-Bandwidth Product (\(GBW\)):}
      We know that,
      \begin{align*}
        GBW &\approx \frac{f_0}{Q} \\
        \text{where } f_0 &= \frac{1}{2\pi \sqrt{L_1 C_1}} \quad \text{(resonant frequency)} \\
      \end{align*}
      \begin{enumerate}
        \item \textbf{Parallel Configuration:}
          \[
            GBW_{parallel} \approx \frac{\frac{1}{2\pi \sqrt{L_1 C_1}}}{\frac{C_1}{G_1} \sqrt{\frac{1}{L_1 C_1}}} = \frac{G_1}{2\pi C_1}
          \]
        \item \textbf{Series Configuration:}
          \[
            GBW_{series} \approx \frac{\frac{1}{2\pi \sqrt{L_1 C_1}}}{G_1 \sqrt{\frac{L_1}{C_1}}} = \frac{1}{2\pi G_1 L_1}
          \]
      \end{enumerate}
    \item \textbf{Impedances:}
      Substituting \(G_1 = \frac{1}{R_1}\) and \(s = j \omega\):
      \begin{enumerate}
        \item \textbf{Parallel Configuration:}
          \begin{enumerate}
            \item At Input:
              \(L_1\) and \(C_1\) are in parallel with each other and
              their combination is in series with \(R_1\):
              \begin{align*}
                Z_{LC} &= \frac{j \omega L_1 \cdot \frac{1}{j \omega C_1}}{j \omega L_1 + \frac{1}{j \omega C_1}} \\
                Z_{LC} &= \frac{L_1}{C_1} \cdot \frac{j \omega C_1}{1 - \omega^2 L_1 C_1} \\
                Z_{LC} &= \frac{L_1 j \omega}{1 - \omega^2 L_1 C_1} \\
                Z_{in} &= R_1 + Z_{LC} \\
                \boxed{Z_{in} = R_1 + \frac{L_1 j \omega}{1 - \omega^2 L_1 C_1}}
              \end{align*}
            \item At Output:
              \(L_1\) and \(C_1\) are in parallel with each other and
              their combination is in parallel with \(R_1\):
              \begin{align*}
                Z_{LC} &= \frac{j \omega L_1 \cdot \frac{1}{j \omega C_1}}{j \omega L_1 + \frac{1}{j \omega C_1}} \\
                Z_{LC} &= \frac{L_1}{C_1} \cdot \frac{j \omega C_1}{1 - \omega^2 L_1 C_1} \\
                Z_{LC} &= \frac{L_1 j \omega}{1 - \omega^2 L_1 C_1} \\
                Z_{out} &= \frac{R_1 \cdot Z_{LC}}{R_1 + Z_{LC}} \\
                Z_{out} &= \frac{R_1 \cdot \frac{L_1 j \omega}{1 - \omega^2 L_1 C_1}}{R_1 + \frac{L_1 j \omega}{1 - \omega^2 L_1 C_1}} \\
                Z_{out} &= \frac{\frac{R_1 L_1 j \omega}{1 - \omega^2 L_1 C_1}}{\frac{R_1(1 - \omega^2 L_1 C_1) + L_1 j \omega}{1 - \omega^2 L_1 C_1}} \\
                Z_{out} &= \frac{R_1 L_1 j \omega}{R_1(1 - \omega^2 L_1 C_1) + L_1 j \omega} \\
                \boxed{Z_{out} = \frac{R_1 L_1 j \omega}{R_1(1 - \omega^2 L_1 C_1) + L_1 j \omega}}
              \end{align*}
          \end{enumerate}
        \item \textbf{Series Configuration:}
          \begin{enumerate}
            \item At Input:
              \(L_1\) and \(C_1\) are in series with each other and
              their combination is in series with \(R_1\):
              \begin{align*}
                Z_{LC} &= j \omega L_1 + \frac{1}{j \omega C_1} \\
                Z_{LC} &= j \left(\omega L_1 - \frac{1}{\omega C_1}\right) \\
                Z_{in} &= R_1 + Z_{LC} \\
                \boxed{Z_{in} = R_1 + j \left(\omega L_1 - \frac{1}{\omega C_1}\right)}
              \end{align*}
            \item At Output:
              \(L_1\) and \(C_1\) are in series with each other and
              their combination is in parallel with \(R_1\):
              \begin{align*}
                Z_{LC} &= j \omega L_1 + \frac{1}{j \omega C_1} \\
                Z_{LC} &= j \left(\omega L_1 - \frac{1}{\omega C_1}\right) \\
                Z_{out} &= \frac{R_1 \cdot Z_{LC}}{R_1 + Z_{LC}} \\
                \boxed{Z_{out} = \frac{R_1 \cdot j \left(\omega L_1 - \frac{1}{\omega C_1}\right)}{R_1 + j \left(\omega L_1 - \frac{1}{\omega C_1}\right)}}
              \end{align*}
            \end{enumerate}
      \end{enumerate}
\end{enumerate}

\subsubsection{Example Design}

\paragraph{Question:}
Design a Band-Pass Filter using both Series and Parallel RLC
Configurations with the following specifications:
\begin{enumerate}
  \item Center Frequency (\(f_0\)): 1 kHz
  \item Quality Factor (\(Q\)): 10
  \item Gain at Center Frequency: Unity (1)
\end{enumerate}
\paragraph{Solution:}
\begin{enumerate}
  \item \textbf{Parallel RLC Configuration:}
    \begin{align*}
      f_0 &= \frac{1}{2\pi \sqrt{L_1 C_1}} = 1 \text{ kHz} \\
      \implies L_1 C_1 &= \frac{1}{(2\pi \cdot 1000)^2} = 2.533 \times 10^{-8} \text{ H-F} \tag{i} \\
      Q &= \frac{C_1}{G_1} \sqrt{\frac{1}{L_1 C_1}} = 10 \\
      \implies G_1 &= \frac{C_1}{10} \sqrt{\frac{1}{L_1 C_1}} = \frac{C_1}{10} \cdot 2\pi \cdot 1000 = 628.32 C_1 \text{ S} \tag{ii}
    \end{align*}
    Assuming \(C_1 = 1 \mu F = 1 \times 10^{-6} F\):
    \begin{align*}
      \text{From (i): } L_1 &= \frac{2.533 \times 10^{-8}}{1 \times 10^{-6}} = 0.02533 H = 25.33 mH \\
      \text{From (ii): } G_1 &= 628.32 \times 1 \times 10^{-6} = 0.00062832 S \\
      \implies R_1 &= \frac{1}{G_1} = \frac{1}{0.00062832} = 1591.55 \Omega
    \end{align*}
    We find the values as:
    \[
      L_1 = 25.33 mH, \quad C_1 = 1 \mu F, \quad R_1 = 1591.55 \Omega
    \]
  \item \textbf{Series RLC Configuration:}
    \begin{align*}
      f_0 &= \frac{1}{2\pi \sqrt{L_1 C_1}} = 1 \text{ kHz} \\
      \implies L_1 C_1 &= \frac{1}{(2\pi \cdot 1000)^2} = 2.533 \times 10^{-8} \text{ H-F} \tag{iii} \\
      Q &= G_1 \sqrt{\frac{L_1}{C_1}} = 10 \\
      \implies G_1 &= 10 \sqrt{\frac{C_1}{L_1}} \tag{iv} \\
    \end{align*}
    Assuming \(C_1 = 1 \mu F = 1 \times 10^{-6} F\):
    \begin{align*}
      \text{From (iii): } L_1 &= \frac{2.533 \times 10^{-8}}{1 \times 10^{-6}} = 0.02533 H = 25.33 mH \\
      \text{From (iv): } G_1 &= 10 \sqrt{\frac{1 \times 10^{-6}}{0.02533}} \approx 0.0628319 S \\
      \implies R_1 &= \frac{1}{G_1} = \frac{1}{0.0628319} \approx 15.9155 \Omega
    \end{align*}
    We find the values as:
    \[
      L_1 = 25.33 mH, \quad C_1 = 1 \mu F, \quad R_1 = 15.9155 \Omega
    \]
\end{enumerate}
